\begin{table}[htbp]
\centering
\footnotesize
\caption{Physical and limnological characteristics of the 11 NTL-LTER primary study lakes \cite{ntllter_lakes434}.}
\label{tab:lake_characteristics}
\resizebox{\textwidth}{!}{%
\begin{tabular}{lccrrll}
\toprule
Lake & Lat (\textdegree N) & Lon (\textdegree W) & Max depth (m) & Area (ha) & Drainage type & Trophic state \\
\midrule
Trout Lake & 46.03 & 89.67 & 35.7 & 1,583 & Drainage & Oligotrophic \\
Big Muskellunge Lake & 46.02 & 89.61 & 21.3 & 376 & Seepage & Oligotrophic \\
Crystal Lake & 46.00 & 89.61 & 20.4 & 38 & Seepage & Oligotrophic \\
Sparkling Lake & 46.01 & 89.70 & 20.0 & 64 & Seepage & Oligotrophic \\
Allequash Lake & 46.04 & 89.62 & 8.0 & 164 & Drainage & Mesotrophic \\
Trout Bog & 46.04 & 89.69 & 7.9 & 1 & Seepage & Dystrophic \\
Crystal Bog & 46.01 & 89.61 & 2.5 & 1 & Seepage & Dystrophic \\
\hline
Lake Mendota & 43.10 & 89.41 & 25.3 & 3,961 & Drainage & Eutrophic \\
Lake Monona & 43.06 & 89.36 & 22.5 & 1,372 & Drainage & Eutrophic \\
Fish Lake & 43.29 & 89.65 & 18.9 & 80 & Seepage & Mesotrophic \\
Lake Wingra & 43.05 & 89.42 & 4.0 & 138 & Drainage & Eutrophic \\
\bottomrule
\end{tabular}%
}
\end{table}
